\chapter{Limitations}
\label{ch:limitations}

% Future Work was promoted from section 10.3 of this chapter to a chapter of
% its own (chapters/11-future-work.tex) on the first examiner's 2026-08-18
% review, registry row R37.

This chapter states what this thesis does not establish. It is written on the principle that a
limitation named by the author is worth more than one found by the reader, and it includes the
hypotheses that were registered before data collection and then never tested --- reported as
such rather than dropped, because a hypothesis that vanishes between registration and reporting
is indistinguishable, to a reader, from one that was tested and inconvenient.

\section{The Larger Corpus Is Synthetic and Born-Digital}
\label{sec:lim-corpus}

The 146-document corpus is generated: its pages are rendered from the same data that supplies
the reference, so they are clean by construction --- no skew, no shadow, no stamp, no staple, no
handwriting, no compression artefact. Real documents in a tax-consulting company arrive with all of those.

The held-out corpus exists precisely to bound this, and it does so with a number:
phone-scanned documents score more than eleven points of mean per-invoice F$_1$ below
email-native ones (Table~\ref{tab:heldout-by-channel}). \emph{Any figure from the synthetic
corpus should therefore be read as an upper bound for the camera-captured case}, and the size of
the correction is \emph{at least} that gap --- the scans here are deskewed scan-app output, the
mildest form of camera capture (Section~\ref{sec:method-data}). The direction is corroborated
independently by a multilingual parsing benchmark reporting an average 17.8\,\% drop for
open-source parsers on photographed documents \parencite{mdpbench2026}.

\section{The Held-Out Corpus Is Small and Single-Annotator}
\label{sec:lim-heldout}

Thirty-nine documents, adjudicated by one person. Two consequences follow.

\textbf{No precise figure is supported.} The headline held-out number should be read as
establishing a level and a direction, not a value to two decimal places. The per-channel cells
are smaller still, some containing around ten documents, so the channel comparison carries the
weight of a clear effect on a small sample, not of a precise estimate.

\textbf{No inter-annotator agreement can be reported.} The three-channel design substantially
limits the exposure --- 463 accepted cells rest on printed-text proof or on two independent
channels agreeing, not on judgement (Section~\ref{sec:method-heldout-gt}) --- but 248 cells were
decided by the author, and for those there is no second opinion. This is irreducible for a
single-author corpus and is the reason the datasheet in Appendix~\ref{app:heldout} records the
provenance class of every cell instead of merely asserting that the key was reviewed.

\section{Header Fields Only, on Real Documents}
\label{sec:lim-groups}

The real-invoice result covers flat header fields. Repeating groups --- line items, the tax
breakdown per rate, discount terms --- are excluded structurally, not by a reporting convention.

The reason is stated in Section~\ref{sec:method-heldout-gt}: rows do not align across reading channels,
because one reader splits a line another merges, so positional comparison manufactures
disagreement out of segmentation rather than out of error. The group rows in the held-out key were
seeded from a single channel and never author-reviewed, and scoring against an unreviewed key is
what produced the retracted figure this project already had to withdraw once.

Line-item extraction is therefore measured on the synthetic corpus, where the reference is exact
and free (Table~\ref{tab:attribution-clusters}). The gap that remains is a measurement of
\emph{real} line items, and Section~\ref{sec:fw-groups} states the route to it.

\section{The Structuring Model Was Held Constant, Not Selected}
\label{sec:lim-structurer-selection}

This is the sharpest asymmetry in the thesis and it deserves to be stated plainly.

The reading model was chosen competitively: a cohort survey, then a shortlist, then a direct
comparison on two instruments, then a hand-judged audit of all 52 residual misses, then a
pre-registered conjunctive test of the larger sibling that the challenger failed
(Section~\ref{sec:results-reader}). \emph{The structuring model was not chosen at all.} It was fixed
as the controlled variable so that the reader comparison and the specification-repair work could
be interpreted against a stable back end, and no structuring bake-off was ever run.

Holding it constant is what makes the rest interpretable --- swapping it would have invalidated
every committed baseline, and the attribution result depends on the structuring stage being the
same in both arms. But the limitation is real and it bounds two claims. No claim can be made that
this is the best available structuring model, and the adaptation results of
Section~\ref{sec:results-adaptation} are conditional on it: a different structuring model might close
the reader-noise gap with no adaptation at all.

There is a related asymmetry worth reporting, because it is the reason the adaptation experiment
was designed as a grid, not a before-and-after. On text rendered from the reference, this
structuring model already scores in the high nineties (Table~\ref{tab:sealed-val-arms}) --- above
the threshold the project had set. The measured shortfall is therefore \emph{robustness to
reading noise}, not capability, and an intervention aimed at capability was always unlikely to
help.

\section{One Adaptation Recipe, No Sweep}
\label{sec:lim-finetune}

The adaptation study tested one recipe: one rank, one learning rate, one schedule, one epoch
budget, 100 fitting examples (Table~\ref{tab:hyperparameters}). No hyperparameter search was
performed, and the training set contained no examples teaching appropriate abstention.

The negative result is therefore about \emph{this} recipe and not about adaptation as a
technique. Section~\ref{sec:disc-finetune} sets out why the observed failure mode --- rising spurious
emission, not degraded placement --- is consistent with published findings, and
Section~\ref{sec:fw-abstention} states the specific untested variant that the mechanism points to.

\section{Hardware Envelope}
\label{sec:lim-hardware}

The hardware claim is two-tier (Section~\ref{sec:method-repro}): a \emph{target class} --- a
workstation with one modern 24\,\ac{GB} graphics processor, the class on which every reading
figure and the entire adaptation study were measured --- and a \emph{floor}, the author's
16\,\ac{GB} laptop, on which the delivered demonstrator runs and the structuring stage is
measured. Both tiers satisfy the privacy premise, which is on-premises processing, and neither
is exotic. The limitation in this section is about how the evidence is distributed across the
two tiers, not about whether the system is local.

The floor bound training and native-resolution reading, not delivered inference. The
adaptation runs exceeded the laptop and were taken on target-class hardware, together with
the precision-matched adaptation evaluations, the reader comparison, and every reported
reading pass; the hardware was rented for the purpose and ran the same open-weights
checkpoints under the author's control (Section~\ref{sec:method-repro}). The local-only property
therefore holds for the delivered pipeline at both tiers --- open-weights end to end, no
network call on any inference path, both models resident within 16\,\ac{GB} --- and a company at
the target class could also adapt the pipeline and reproduce the full measurement record on
premises. A company equipped only at the floor could run the pipeline but not adapt it or
reproduce the measurements.

\textbf{The sharper form of this limitation, stated plainly.} On the laptop the reader's input
is capped to roughly half an A4 page's resolution --- a Metal-memory workaround specific to
this machine, not a quality choice and not a property of local deployment, since the same
reader handles native resolution within a single 24\,\ac{GB} card. Every transcript in this
thesis was produced above that cap, at the target class. The reading quality of the floor
configuration was therefore not measured, on either corpus, and no figure reported here should
be read as a floor figure.

This matters more than a routine venue note, because it lands on the axis this thesis identifies
as the dominant real-world cost. The channel comparison of Section~\ref{sec:results-heldout} shows
capture quality costing more than eleven points of mean per-invoice F$_1$, and page resolution
is a capture-quality variable. The unmeasured floor penalty and the measured channel penalty are
plausibly the same mechanism, which means the floor figure should be expected to sit below the
reported one --- by an amount this thesis cannot state. Both directions of error are available to
a reader who wants to speculate; the honest position is that the number is missing, not small.

The measurement that would close it is bounded, requires no rented hardware and no new code,
and is specified in Section~\ref{sec:fw-envelope}. It was not taken within the scope of this thesis
because the experiment track was frozen in favour of the writeup
(Section~\ref{sec:design-scope}), and recording the gap was judged better than an unplanned run
reported without the pre-registration discipline the rest of the measurements carry.

\section{Language, Document Class and Jurisdiction}
\label{sec:lim-scope}

German and English business-to-business invoices, under German and European invoicing
conventions. Nothing is claimed about receipts, contracts, bank statements, dunning letters or
correspondence; nothing about jurisdictions with different mandatory particulars; nothing about
languages outside the two measured. The single French-language document in the synthetic corpus is
an instance, not evidence of transfer.

\section{Measurement Constructs}
\label{sec:lim-measurement}

Two constructs carry more inferential weight than their definitions strictly support.

\textbf{The reading-quality proxy is a proxy.} It asks only whether a reference value appears
somewhere in the transcript. A transcript that jumbled every value on the page would satisfy it
perfectly and still be useless downstream (Section~\ref{sec:method-answerability}). It was
consequently never the sole basis for a decision, and the reader decision it touched was
settled by adjudicating the individual misses.

\textbf{F$_1$ over header fields is not a measure of legal usability.} It measures extraction
correctness. It does not measure whether a booking would survive audit, whether the mandatory
particulars of \S~14~\ac{UStG} are satisfied, or whether an input-tax deduction is defensible. The
weighted compliance metric that would have gestured toward this was abandoned deliberately
(Section~\ref{sec:method-metrics}), and no legal assessment is offered anywhere in this thesis.

\section{Three Known Defects Remain Open}
\label{sec:lim-open-defects}

Reported here so that the apparatus is not represented as cleaner than it is.

\textbf{Twenty-two answer-key cells are withdrawn from scoring on no criterion.} The scoring
apparatus uses one internal marker for two unrelated things: a deliberate, rule-driven
neutralisation, and a signed-off value the parser could not read. The second is a locale-parsing
gap, and because a withdrawn cell leaves the numerator and denominator alike, no reported figure
could have revealed it (Section~\ref{sec:validity-open}). The affected cells are all dates and
concentrate in the capture channel the system reads best, so the reported figure is more likely
understated than flattered --- but the honest position is that the direction is undetermined
until the parser is repaired and the corpus re-scored on frozen generations.

\textbf{One field regressed under a correct repair and stays regressed.} Making the
perfect-text renderer faithful to the reference exposed a comparison problem in the supplier
article number field that was not resolved (Section~\ref{sec:validity-open}).

\textbf{A stale configuration default remains a hazard.} A diagnostic that constructs the
experiment configuration from its defaults, instead of loading the project's configuration file,
silently selects a superseded reading model. The two audit tools were fixed to load the
configuration explicitly; the default itself was not changed, so any future script that builds
the configuration bare will measure the wrong pipeline (Section~\ref{sec:validity-specification}).

\section{The Evaluation Methodology Is Not Itself Deployable}
\label{sec:lim-cloud-gt}

Building the held-out answer key required two cloud services, so the content of those 39 invoices
left local hardware (Section~\ref{sec:method-heldout-gt}). The distinction relied upon --- a measuring
instrument is not part of the delivered system --- holds for the pipeline, which makes no network
call on any inference path. It does not hold for the methodology: a practitioner bound by
professional secrecy could not construct this answer key from their own client documents.
Section~\ref{sec:fw-local-gt} states what closing that would require.

\section{Pre-Registered but Unevaluated Hypotheses}
\label{sec:unevaluated}

The scope of this thesis was narrowed during the work, in favour of one complete and honestly
measured layer over three partially measured ones. Five registered hypotheses were consequently
never tested. They are listed here as \emph{not evaluated}, which is distinct from refuted and
distinct from inconclusive: no data was collected against them, so they remain open. A sixth item
is included because it is routinely mistaken for one of them, and was not registered at all.

\begin{enumerate}
  \item \textbf{Knowledge-graph fidelity.} That extracted records could be assembled into a
        graph whose entity resolution is accurate enough for cross-document reasoning. The layer
        was designed (Section~\ref{sec:design-layer2}) and not built.
  \item \textbf{Graph versus flat retrieval.} That graph-structured retrieval would outperform
        flat vector retrieval on multi-document analytical questions. Untested: it presupposes
        the graph.
  \item \textbf{Query routing.} That a router could reliably dispatch a natural-language question
        to the appropriate retrieval strategy. Untested for the same reason.
  \item \textbf{Conditional validator retry.} That re-prompting the structuring stage with the
        specific validation failure it produced would recover a measurable share of failures. The
        validate-and-repair stage was built deterministic and single-pass; the retry variant was
        never implemented, and no figure in this thesis reflects one.
  \item \textbf{Cloud versus local comparison.} That the accuracy cost of the local-only
        constraint could be quantified against a frontier cloud model on the same corpus.
        Deliberately not run on client documents, for the reason the thesis exists;
        Section~\ref{sec:fw-cloud} describes the form it could safely take.
  \item \textbf{Template shift --- floated, never registered.} That performance would degrade
        measurably on invoice layouts drawn from vendors absent from the fitting set. This was
        discussed as a candidate and never entered in the register
        (Appendix~\ref{app:hypotheses}), so no pre-registered claim attaches to it. It was
        partially and accidentally addressed by the held-out corpus, whose layouts were all
        unseen --- but not in the controlled form that had been contemplated, since layout
        novelty and capture channel vary together there and cannot be separated.
\end{enumerate}
